\chapter{Composition Errors, Decoded}
\label{app:comperrors}

These entries refer to the pinned composer in appendix~\ref{app:toolchain}.
A schema coordinate such as \code{Session.averageScore} names a type and
field, not a source-file line. The quoted diagnostic fragments below are
excerpts of the full captures in chapters~\ref{ch:composition}
and~\ref{ch:satisfiability}. The stage matters: a compiler error, a composition
error and a router execution error require different evidence.

\section{Directive Validation}

\subsection*{Cost Weight Has the Wrong Type}

Symptom: \code{@cost(weight: ...)} reports a value that is not a valid
\code{Int!}. In this example the rejected value is a quoted string, so the
diagnostic contains doubled quotation marks.

Cause: Hot Chocolate's generated cost directive and Cosmo's definition
disagree. This can fail with only one subgraph; no ownership collision is
needed.

Action: use the complete \code{Program.cs} in
section~\ref{sec:ch07-the-second-document}, which disables automatic cost
defaults with \code{ModifyCostOptions}. Rebuild and export before composing.
Do not hand-edit every generated weight and leave the next export to undo it.

\subsection*{List-Size Argument Is Unknown}

Symptom: the \code{@listSize} diagnostic says its definition does not define
an argument that was provided. The captured console box cuts off
\code{slicingArgumentDefaultValue}.

Cause: the two tools disagree about that directive's arguments.

Action: inspect the untruncated exported SDL to identify the argument.
The same automatic-cost-default change removes the conflicting annotation
in this example. Disabling resolvability validation does not disable
directive validation.

\subsection*{Override Label Is Rejected}

Symptom: the composer rejects the \code{label} argument on \code{@override}.

Cause: the Hot Chocolate attribute exposes progressive override, but the
pinned Cosmo composer does not accept that argument.

Action: do not infer cross-product support from a compiling attribute.
Use the non-progressive migration described in
section~\ref{sec:ch17-handing-a-field-over}, or verify a different toolchain
as a separate change.

\section{Conflicting Declarations}

\subsection*{Field Is Declared More Than Once}

Symptom excerpt:

\begin{minted}[fontsize=\footnotesize,breakanywhere]{text}
The Object "Speaker" defines the same fields in multiple subgraphs without the "@shareable" directive:
\end{minted}

Cause: more than one subgraph claims the same field without declarations
that permit sharing.

Action: locate every named declaration. If one service is only pointing at
a speaker, reduce it to the key-only stub. If the field is genuinely
shareable, make the required declarations agree and test equivalent answers.
Adding a directive solely to silence the error does not establish equivalence.

\subsection*{Generated Node Fields Collide}

Symptom: the same error names \code{Query.node} or \code{Query.nodes},
although no application method declared either field.

Cause: global object identification generated them in more than one service.

Action: the \code{ShareNodeFields} interceptor in
section~\ref{sec:ch07-the-second-document} targets those two fields only.
Do not mark the whole \code{Query} type shareable to solve this collision.
Chapter~\ref{ch:second-third-subgraph} documents the remaining node-routing
limitation even after composition accepts the declarations.

\subsection*{Two Output Types Merge Without an Error}

Symptom: one subgraph declares \code{[Node]!}, another \code{[Node]}, and
composition succeeds with the nullable result.

Cause: the composed output accepts the more permissive value. This is not
the earlier directive-validation failure.

Action: compare the composed field and each direct subgraph contract.
Chapter~\ref{ch:null-blast-radius} makes both services nullable to align
their direct and routed failure boundaries; a successful merge alone does
not make those boundaries identical.

\section{A Field Has No Reachable Entity Route}

\subsection*{Unresolvable Field}

Symptom excerpt:

\begin{minted}{text}
The field "averageScore" is unresolvable at the following path:
\end{minted}

Cause: the composer cannot reach the named field along the reported path.
In the book's Ratings experiment, \code{Session} carries a non-resolvable
key even though Ratings must accept entity lookups for its contributed
fields.

Action: read the reason naming the field's subgraph, then inspect its key,
key fields and reference resolver. Check whether the originating service can
supply that key. The printed root is one route that failed, not necessarily
the service to edit. See section~\ref{sec:ch15-reading-one-properly}.

\subsection*{Interface Has No Object Implementation}

Symptom excerpt:

\begin{minted}[fontsize=\footnotesize,breakanywhere]{text}
Subgraph "speakers": The Interface "Node" is used as an output type without at least one Object type implementation
\end{minted}

Cause: the experiment kept the node root while removing the implementing
object from that document.

Action: distinguish this warning from the unreachable-field defect.
The pinned composer can emit it and still accept the broken graph.
Restore the normal exported document before continuing; the warning is not
a repair.

\subsection*{Disabling the Check Produces a Config}

Symptom: composition succeeds only when the resolvability check is disabled
with this flag:

\begin{minted}{text}
--disable-resolvability-validation
\end{minted}

Cause: the reachability pass was skipped. Structural and directive checks
still ran.

Action: use the flag only to isolate which pass raised the diagnostic.
A config written under that flag is not evidence that its fields can be
served. The runnable example does not retain the flag.

\section{Composition Passes, Requests Fail}

\subsection*{The False Pass in the Three-Service Variant}

Symptom: the broken Ratings key is accepted when Search is absent and the
shared singular node field is visited first.

Cause: the pinned composer's visited-entity bookkeeping skips paths it has
not actually proved resolvable. Reordering root fields changes detection
without repairing the key.

Action: use the named false-pass input in
section~\ref{sec:ch15-the-composition-that-passes}. Compare the reduced
Sessions/Ratings pair and the reordered-root control, then send requests for
every contributed field. Repeating an identical successful composition adds
no coverage.

\subsection*{Three Runtime Outcomes of That False Pass}

\begin{center}
\small
\begin{tabular}{p{0.25\linewidth}p{0.62\linewidth}}
\toprule
Selected field & Observed result \tabularnewline
\midrule
\code{averageScore} & Router planning fails, HTTP 500 \tabularnewline
\code{ratingCount} & Wrong root-field fetch; client positions null with an error \tabularnewline
\code{feedbackUrl} & Null positions with no error entry \tabularnewline
\bottomrule
\end{tabular}
\end{center}

These are the captured results of one pinned defect, not general GraphQL
rules. Nullable positions limit propagation; they do not require errors to
be omitted. Restore the resolvable Ratings key, its reference resolver and
the four-service configuration after the experiment.

\subsection*{Client Compatibility or Ownership Is Wrong}

Symptom: composition succeeds after a public field is removed, an input is
made required, or an existing field is rerouted to another service.

Cause: composition is not a comparison with the published client contract,
a retained operation corpus or a record of team approval.

Action: compare the client-facing composed schemas, classify compatibility
changes, replay retained operations and review routing changes with the
responsible team. Chapters~\ref{ch:breaking-changes}
and~\ref{ch:ownership} separate those checks from composition.
